\documentclass[11pt]{article}

\usepackage[a4paper,margin=1in]{geometry}
\usepackage[T1]{fontenc}
\usepackage[utf8]{inputenc}
\usepackage{amsmath,amssymb,amsthm,mathtools}
\usepackage{enumitem}
\usepackage{microtype}
\usepackage[hidelinks]{hyperref}

\setlist[itemize]{leftmargin=2em}
\setlist[enumerate]{leftmargin=2em}

\newcommand{\F}{\mathbb F}
\newcommand{\R}{\mathbb R}
\newcommand{\Z}{\mathbb Z}
\newcommand{\RP}{\mathbb{RP}}
\newcommand{\Gr}{\operatorname{Gr}}
\newcommand{\id}{\operatorname{id}}
\newcommand{\rank}{\operatorname{rank}}
\newcommand{\Hom}{\operatorname{Hom}}

\theoremstyle{plain}
\newtheorem{theorem}{Theorem}[section]
\newtheorem{proposition}[theorem]{Proposition}
\newtheorem{lemma}[theorem]{Lemma}
\newtheorem{corollary}[theorem]{Corollary}

\theoremstyle{definition}
\newtheorem{definition}[theorem]{Definition}
\newtheorem{remark}[theorem]{Remark}

\numberwithin{equation}{section}

\title{Computing the Ring of $\mathbb Z/2\mathbb Z$-Characteristic Classes\\
\large Submission Note}
\author{W. Huang}
\date{12 June 2026}

\begin{document}

\maketitle

\begin{abstract}
This note computes the mod $2$ cohomology ring of the infinite real Grassmannian
\[
  G_n=\Gr_n(\R^\infty),
\]
and explains how this computation implies the uniqueness of Stiefel--Whitney classes.  Let
\[
  \gamma_n\longrightarrow G_n
\]
be the tautological rank $n$ real vector bundle.  Assuming the usual axioms for Stiefel--Whitney classes, we prove
\[
  H^*(G_n;\F_2)\cong \F_2[w_1,\ldots,w_n],
  \qquad |w_i|=i,
\]
where $w_i=w_i(\gamma_n)$.  The proof has two parts.  First, one pulls the tautological bundle back to $(\RP^\infty)^n$ and uses the resulting split bundle to prove that $w_1,\ldots,w_n$ are algebraically independent.  Second, one compares the number of monomials in the $w_i$'s with the number of Schubert cells in $G_n$.  Finally, the same pullback argument proves that any two systems of characteristic classes satisfying the Stiefel--Whitney axioms must agree.
\end{abstract}

\tableofcontents

\section{Introduction}

Characteristic classes are cohomology classes naturally associated to vector bundles.  They provide cohomological invariants of the way a vector bundle is twisted over its base space.  For a real vector bundle
\[
  \xi\longrightarrow B,
\]
the Stiefel--Whitney classes are classes
\[
  w_i(\xi)\in H^i(B;\F_2),
\]
and it is convenient to collect them into the total class
\[
  w(\xi)=1+w_1(\xi)+w_2(\xi)+\cdots.
\]

The aim of this note is not to construct the Stiefel--Whitney classes from scratch.  Instead, we assume that a system of classes satisfying the Stiefel--Whitney axioms exists, and we study what those axioms force.  The main result is the universal computation
\[
  H^*(\Gr_n(\R^\infty);\F_2)
  \cong
  \F_2[w_1,\ldots,w_n].
\]
This is a precise form of the statement that every mod $2$ characteristic class of rank $n$ real vector bundles is a polynomial in the Stiefel--Whitney classes.

The proof follows the standard strategy of Milnor--Stasheff, Section 7.  The first step is a splitting test: the universal bundle is pulled back to the product $(\RP^\infty)^n$, where it becomes a direct sum of line bundles.  This shows that the universal classes pull back to elementary symmetric polynomials, hence satisfy no algebraic relation.  The second step uses the Schubert cell decomposition of the infinite Grassmannian.  The number of Schubert cells gives an upper bound for cohomology in each degree, while the algebraically independent monomials in the $w_i$'s give a matching lower bound.  Therefore the monomials in the $w_i$'s form a basis in every degree.

After proving the ring computation, we prove uniqueness.  If $w$ and $\overline w$ are two possible assignments of total Stiefel--Whitney classes satisfying the axioms, then they agree first on line bundles, then on split bundles over $(\RP^\infty)^n$, then on the universal bundle over $G_n$, and finally on all vector bundles by pullback along classifying maps.

\section{Universal bundles and characteristic classes}

\subsection{The infinite Grassmannian}

Throughout the note we write
\[
  \F_2=\Z/2\Z,
  \qquad
  G_n=\Gr_n(\R^\infty).
\]
The space $G_n$ is the infinite real Grassmannian: its points are $n$-dimensional linear subspaces of $\R^\infty$.  It carries the tautological rank $n$ real vector bundle
\[
  \gamma_n\longrightarrow G_n.
\]
The fiber of $\gamma_n$ over a point $V\in G_n$ is the vector space $V$ itself.

The fundamental reason that $G_n$ appears is the classification theorem for vector bundles.

\begin{theorem}[Classification theorem]
Let $B$ be a paracompact space.  Then isomorphism classes of rank $n$ real vector bundles over $B$ are naturally classified by homotopy classes of maps
\[
  B\longrightarrow G_n.
\]
More precisely, if $\xi\to B$ is a rank $n$ real vector bundle, then there exists a map
\[
  f:B\to G_n
\]
such that
\[
  \xi\cong f^*\gamma_n.
\]
The classifying map $f$ is unique up to homotopy.
\end{theorem}

The paracompactness hypothesis is important in the usual classification theorem.  In this note all base spaces to which we apply the theorem, such as CW complexes and finite products of countable CW complexes, are paracompact.

\subsection{The ring of characteristic classes}

For fixed rank $n$, a mod $2$ characteristic class of degree $k$ is a natural rule assigning to every rank $n$ real vector bundle $\xi\to B$ a class
\[
  c(\xi)\in H^k(B;\F_2).
\]
Naturality means that for every map $h:B'\to B$ one has
\[
  c(h^*\xi)=h^*c(\xi).
\]

The classification theorem identifies such characteristic classes with cohomology classes of $G_n$.  Indeed, if
\[
  c\in H^k(G_n;\F_2)
\]
and if $f:B\to G_n$ classifies $\xi$, then one defines
\[
  c(\xi)=f^*c.
\]
This is independent of the choice of $f$, because homotopic maps induce the same map in cohomology.  Conversely, any natural characteristic class is determined by its value on the universal bundle $\gamma_n$.  Thus the ring of mod $2$ characteristic classes for rank $n$ real vector bundles is naturally
\[
  H^*(G_n;\F_2).
\]

Under this interpretation, computing $H^*(G_n;\F_2)$ is the same as computing all mod $2$ characteristic classes of rank $n$ real vector bundles.

\subsection{Stiefel--Whitney axioms used here}

We will use the following standard properties of Stiefel--Whitney classes.

\begin{enumerate}
  \item \textbf{Degree and rank.}  If $\xi\to B$ has rank $r$, then
  \[
    w_i(\xi)\in H^i(B;\F_2),
    \qquad
    w_0(\xi)=1,
    \qquad
    w_i(\xi)=0 \text{ for } i>r.
  \]

  \item \textbf{Naturality.}  If $f:B'\to B$, then
  \[
    w_i(f^*\xi)=f^*w_i(\xi).
  \]

  \item \textbf{Whitney product formula.}  If $\xi$ and $\eta$ are real vector bundles over the same base, then
  \[
    w(\xi\oplus\eta)=w(\xi)w(\eta).
  \]

  \item \textbf{Normalization.}  If $\gamma_1\to\RP^1$ is the tautological line bundle and $a\in H^1(\RP^1;\F_2)$ is the nonzero class, then
  \[
    w(\gamma_1)=1+a.
  \]
\end{enumerate}

The proof below mostly uses the tautological line bundle over $\RP^\infty$, still denoted by
\[
  \gamma_1\longrightarrow \RP^\infty.
\]
The normalization over $\RP^1$ implies
\[
  w(\gamma_1)=1+a
\]
for the tautological bundle over $\RP^\infty$, where
\[
  a\in H^1(\RP^\infty;\F_2)
\]
is the standard generator.  This follows because the restriction map
\[
  H^1(\RP^\infty;\F_2)	o H^1(\RP^1;\F_2)
\]
is an isomorphism, and a line bundle has no Stiefel--Whitney classes above degree $1$.

\subsection{The cohomology of $(\RP^\infty)^n$}

We shall use the standard calculation
\[
  H^*(\RP^\infty;\F_2)
  \cong
  \F_2[a],
  \qquad |a|=1.
\]
Here is the calculation.  For finite $m$, the usual CW structure on $\RP^m$ has one cell in each dimension $0,1,\ldots,m$, and with $\F_2$-coefficients the cellular differential is zero.  The standard cup-product computation gives
\[
  H^*(\RP^m;\F_2)
  \cong
  \F_2[a]/(a^{m+1}),
  \qquad |a|=1.
\]
Passing to the limit over the inclusions $\RP^m\subset\RP^\infty$ removes the relation $a^{m+1}=0$, so
\[
  H^*(\RP^\infty;\F_2)
  \cong
  \F_2[a].
\]
This is the projective-space calculation used in the proof of the universal Stiefel--Whitney class theorem; see Milnor--Stasheff~\cite[Section 7]{MS74}.  For the cellular and cup-product computation of projective spaces, see also Hatcher~\cite[Section 3.2]{Hat02}.

Now let
\[
  X=(\RP^\infty)^n,
\]
and let
\[
  \pi_i:X\to\RP^\infty
\]
be the projection onto the $i$-th factor.  Define
\[
  a_i=\pi_i^*a\in H^1(X;\F_2).
\]
Since the coefficient ring is the field $\F_2$, the Kunneth theorem has no Tor terms and gives
\[
  H^*(X;\F_2)
  \cong
  H^*(\RP^\infty;\F_2)^{\otimes n}
  \cong
  \F_2[a_1,\ldots,a_n].
\]
This polynomial ring will be the main testing ground for the universal Stiefel--Whitney classes.

\section{The main theorem}

\begin{theorem}\label{thm:main}
Let
\[
  G_n=\Gr_n(\R^\infty),
\]
and let
\[
  \gamma_n\to G_n
\]
be the tautological rank $n$ real vector bundle.  Then
\[
  H^*(G_n;\F_2)
  \cong
  \F_2[w_1,\ldots,w_n],
  \qquad |w_i|=i,
\]
where
\[
  w_i=w_i(\gamma_n).
\]
Equivalently, the graded ring homomorphism
\[
  \Phi:\F_2[x_1,\ldots,x_n]
  \longrightarrow
  H^*(G_n;\F_2),
  \qquad
  x_i\longmapsto w_i(\gamma_n),
\]
with $|x_i|=i$, is an isomorphism.
\end{theorem}

The proof is divided into two steps.  First we prove that $\Phi$ is injective, or equivalently that $w_1,\ldots,w_n$ satisfy no nontrivial polynomial relation.  Then we prove that these classes generate all of $H^*(G_n;\F_2)$ by a dimension comparison using the Schubert cell decomposition.

\section{Algebraic independence}

Let
\[
  X=(\RP^\infty)^n.
\]
For each $i$, define the line bundle
\[
  \ell_i=\pi_i^*\gamma_1
\]
over $X$.  Consider the rank $n$ bundle
\[
  \xi=\ell_1\oplus\cdots\oplus\ell_n.
\]
By naturality,
\[
  w(\ell_i)
  =w(\pi_i^*\gamma_1)
  =\pi_i^*w(\gamma_1)
  =1+a_i.
\]
Therefore the Whitney product formula gives
\[
  w(\xi)
  =\prod_{i=1}^n w(\ell_i)
  =\prod_{i=1}^n(1+a_i).
\]
Expanding this product yields
\[
  w(\xi)
  =1+e_1(a_1,\ldots,a_n)+e_2(a_1,\ldots,a_n)+\cdots+e_n(a_1,\ldots,a_n),
\]
where
\[
  e_k(a_1,\ldots,a_n)
  =
  \sum_{1\leq i_1<\cdots<i_k\leq n}
  a_{i_1}\cdots a_{i_k}
\]
is the $k$-th elementary symmetric polynomial.  Thus
\[
  w_k(\xi)=e_k(a_1,\ldots,a_n).
\]

The bundle $\xi$ is a rank $n$ real vector bundle over the paracompact space $X$.  Hence it is classified by a map
\[
  g:X\to G_n
\]
such that
\[
  \xi\cong g^*\gamma_n.
\]
Naturality gives
\[
  g^*w_k(\gamma_n)
  =w_k(g^*\gamma_n)
  =w_k(\xi)
  =e_k(a_1,\ldots,a_n).
\]
Therefore the map
\[
  g^*:H^*(G_n;\F_2)	o H^*(X;\F_2)=\F_2[a_1,\ldots,a_n]
\]
sends the universal Stiefel--Whitney class $w_k(\gamma_n)$ to the elementary symmetric polynomial $e_k$.

\begin{lemma}\label{lem:algindep}
The elementary symmetric polynomials
\[
  e_1(a_1,\ldots,a_n),\ldots,e_n(a_1,\ldots,a_n)
\]
are algebraically independent over $\F_2$.
\end{lemma}

\begin{proof}
By the fundamental theorem of symmetric polynomials,
\[
  \F_2[a_1,\ldots,a_n]^{S_n}
  =
  \F_2[e_1,\ldots,e_n],
\]
and every symmetric polynomial has a unique expression as a polynomial in the elementary symmetric polynomials.  If there were a nonzero polynomial
\[
  P\in \F_2[x_1,\ldots,x_n]
\]
with
\[
  P(e_1,\ldots,e_n)=0,
\]
then the zero symmetric polynomial would have two distinct expressions in terms of the $e_i$'s.  This contradicts uniqueness.  Hence the $e_i$'s are algebraically independent.
\end{proof}

Now suppose that
\[
  P(w_1,\ldots,w_n)=0
\]
for some polynomial $P\in\F_2[x_1,\ldots,x_n]$.  Pulling back along $g$ gives
\[
  0
  =g^*P(w_1,\ldots,w_n)
  =P(g^*w_1,\ldots,g^*w_n)
  =P(e_1,\ldots,e_n).
\]
By Lemma~\ref{lem:algindep}, this implies $P=0$.  Thus the map
\[
  \Phi:\F_2[x_1,\ldots,x_n]\to H^*(G_n;\F_2)
\]
is injective.

\section{Schubert-cell counting}

It remains to show that the classes $w_1,\ldots,w_n$ generate all of $H^*(G_n;\F_2)$.  This is where the Schubert cell structure of $G_n$ enters.

The Schubert cells of $G_n$ are indexed by partitions
\[
  \lambda=(\lambda_1,\ldots,\lambda_n)
\]
with
\[
  \lambda_1\geq \lambda_2\geq\cdots\geq \lambda_n\geq 0.
\]
The cell corresponding to $\lambda$ has dimension
\[
  |\lambda|=\lambda_1+\cdots+\lambda_n.
\]
Therefore the number of cells of dimension $r$ is the number of partitions of $r$ into at most $n$ parts.  We denote this number by
\[
  p_{\leq n}(r).
\]

Cellular cochains give an upper bound:
\[
  \dim_{\F_2}H^r(G_n;\F_2)
  \leq
  \dim_{\F_2}C^r(G_n;\F_2)
  =
  p_{\leq n}(r).
\]
This is only an upper bound.  We do not need to compute the cellular differential, because the algebraically independent Stiefel--Whitney monomials will give the matching lower bound.

Indeed, since $|w_i|=i$, a monomial
\[
  w_1^{r_1}w_2^{r_2}\cdots w_n^{r_n}
\]
has degree
\[
  r_1+2r_2+\cdots+nr_n.
\]
Thus the number of degree $r$ monomials in $w_1,\ldots,w_n$ is the number of nonnegative integer solutions of
\[
  r_1+2r_2+\cdots+nr_n=r.
\]
Equivalently, it is the number of partitions of $r$ whose parts have size at most $n$.  By conjugation of Young diagrams, this number is equal to the number of partitions of $r$ into at most $n$ parts, namely $p_{\leq n}(r)$.

Since $\Phi$ is injective, the degree $r$ monomials in $w_1,\ldots,w_n$ are linearly independent in $H^r(G_n;\F_2)$.  There are exactly $p_{\leq n}(r)$ such monomials, so
\[
  \dim_{\F_2}H^r(G_n;\F_2)
  \geq
  p_{\leq n}(r).
\]
Together with the cellular upper bound, this gives
\[
  \dim_{\F_2}H^r(G_n;\F_2)
  =
  p_{\leq n}(r)
\]
for every $r$.  Therefore the degree $r$ monomials in $w_1,\ldots,w_n$ form a basis of $H^r(G_n;\F_2)$ for every $r$.  Hence $\Phi$ is surjective as well as injective, and Theorem~\ref{thm:main} follows.

\section{Examples}

\subsection{The case $n=1$}

For $n=1$,
\[
  G_1=\Gr_1(\R^\infty)=\RP^\infty.
\]
The theorem gives
\[
  H^*(G_1;\F_2)=\F_2[w_1].
\]
This is exactly the standard computation
\[
  H^*(\RP^\infty;\F_2)=\F_2[a],
\]
with
\[
  w_1(\gamma_1)=a.
\]
Thus the case $n=1$ recovers the normalization for the tautological line bundle.

\subsection{The case $n=2$}

For $n=2$,
\[
  H^*(G_2;\F_2)=\F_2[w_1,w_2],
  \qquad |w_1|=1,
  \quad |w_2|=2.
\]
For example, the first few graded pieces have bases
\[
  H^0:\quad 1,
\]
\[
  H^1:\quad w_1,
\]
\[
  H^2:\quad w_1^2,\;w_2,
\]
\[
  H^3:\quad w_1^3,\;w_1w_2,
\]
and
\[
  H^4:\quad w_1^4,\;w_1^2w_2,\;w_2^2.
\]
This illustrates what the polynomial-ring theorem means in practice: there are no relations between $w_1$ and $w_2$.

\subsection{The case $n=3$}

For $n=3$,
\[
  H^*(G_3;\F_2)=\F_2[w_1,w_2,w_3].
\]
In degree $4$, a basis is
\[
  w_1^4,
  \qquad
  w_1^2w_2,
  \qquad
  w_2^2,
  \qquad
  w_1w_3.
\]
These four monomials correspond to the four partitions of $4$ into at most three parts:
\[
  4,
  \qquad
  3+1,
  \qquad
  2+2,
  \qquad
  2+1+1.
\]

\section{Uniqueness of Stiefel--Whitney classes}

We now use the computation of $H^*(G_n;\F_2)$ to prove uniqueness.  Suppose that
\[
  \xi\longmapsto w(\xi)
\]
and
\[
  \xi\longmapsto \overline w(\xi)
\]
are two assignments of total Stiefel--Whitney classes to real vector bundles, and suppose that both assignments satisfy the Stiefel--Whitney axioms.  We will prove that
\[
  w_i(\xi)=\overline w_i(\xi)
\]
for every real vector bundle $\xi$ and every $i\geq 0$.

\subsection{Agreement on line bundles}

First consider the tautological line bundle
\[
  \gamma_1\to\RP^\infty.
\]
The normalization axiom, together with the restriction to $\RP^1$, implies that both assignments give
\[
  w(\gamma_1)=1+a,
  \qquad
  \overline w(\gamma_1)=1+a.
\]
Thus the two assignments agree on the universal line bundle over $\RP^\infty$.  Since every real line bundle over a paracompact base is pulled back from this universal line bundle, the two assignments agree on all real line bundles.

\subsection{Agreement on the universal split bundle}

Let
\[
  X=(\RP^\infty)^n
\]
and consider the split bundle
\[
  \xi=\ell_1\oplus\cdots\oplus\ell_n,
  \qquad
  \ell_i=\pi_i^*\gamma_1.
\]
Since the two assignments agree on each $\ell_i$, the Whitney product formula gives
\[
  w(\xi)
  =
  \prod_{i=1}^n(1+a_i)
  =
  \overline w(\xi).
\]
Thus the two assignments agree on this split rank $n$ bundle over $X$.

\subsection{Agreement on the universal rank $n$ bundle}

Let
\[
  g:X\to G_n
\]
classify the bundle $\xi$, so that
\[
  \xi\cong g^*\gamma_n.
\]
By naturality,
\[
  g^*w(\gamma_n)=w(\xi),
  \qquad
  g^*\overline w(\gamma_n)=\overline w(\xi).
\]
Since $w(\xi)=\overline w(\xi)$, we have
\[
  g^*w(\gamma_n)=g^*\overline w(\gamma_n).
\]

We now use the ring computation.  Let
\[
  u_i=w_i(\gamma_n)
\]
denote the universal classes coming from the first assignment $w$.  By Theorem~\ref{thm:main},
\[
  H^*(G_n;\F_2)=\F_2[u_1,\ldots,u_n].
\]
Moreover,
\[
  g^*u_i=e_i(a_1,\ldots,a_n).
\]
Since the elementary symmetric polynomials are algebraically independent, the induced homomorphism
\[
  g^*:H^*(G_n;\F_2)	o H^*(X;\F_2)
\]
is injective.  Hence the equality after pulling back by $g$ implies
\[
  w(\gamma_n)=\overline w(\gamma_n).
\]
Therefore
\[
  w_i(\gamma_n)=\overline w_i(\gamma_n)
\]
for all $i$.

\begin{remark}
The notation $u_i$ is used only to avoid confusing the two assignments.  The classes $\overline w_i(\gamma_n)$ also lie in $H^*(G_n;\F_2)$, and after the main theorem every element of this cohomology ring is a polynomial in $u_1,\ldots,u_n$.
\end{remark}

\subsection{Agreement on arbitrary bundles}

Let
\[
  \eta\to B
\]
be any rank $n$ real vector bundle over a paracompact space $B$.  Choose a classifying map
\[
  f:B\to G_n
\]
with
\[
  \eta\cong f^*\gamma_n.
\]
Using naturality and the equality on the universal bundle, we get
\[
  w(\eta)
  =f^*w(\gamma_n)
  =f^*\overline w(\gamma_n)
  =\overline w(\eta).
\]
Thus the two assignments agree on every rank $n$ real vector bundle.  Since $n$ is arbitrary, they agree on all real vector bundles.  This proves the uniqueness of Stiefel--Whitney classes.

\section{Conclusion}

The main result of this note is the ring computation
\[
  H^*(\Gr_n(\R^\infty);\F_2)
  \cong
  \F_2[w_1,\ldots,w_n],
  \qquad |w_i|=i.
\]
The proof combines two complementary ideas.  Pulling back to $(\RP^\infty)^n$ gives a split bundle whose Stiefel--Whitney classes are elementary symmetric polynomials.  This proves that the universal classes $w_1,\ldots,w_n$ are algebraically independent.  The Schubert cell decomposition then shows that in each degree there are no more cohomology classes than the monomials in the $w_i$'s.  Therefore the Stiefel--Whitney monomials form an additive basis, and the cohomology ring is a polynomial ring.

As an application, the same computation proves uniqueness.  The axioms force any Stiefel--Whitney assignment to agree with the standard one on line bundles, then on split bundles, then on the universal bundle, and finally on every vector bundle by naturality and classification.

\appendix

\section{A direct proof of algebraic independence}

For completeness, here is a direct proof of Lemma~\ref{lem:algindep}.  Use a lexicographic monomial order with
\[
  a_1>a_2>\cdots>a_n.
\]
The leading term of $e_k(a_1,\ldots,a_n)$ is
\[
  a_1a_2\cdots a_k.
\]
Therefore the leading term of
\[
  e_1^{m_1}e_2^{m_2}\cdots e_n^{m_n}
\]
is
\[
  a_1^{m_1+\cdots+m_n}
  a_2^{m_2+\cdots+m_n}
  \cdots
  a_n^{m_n}.
\]
The exponent vector uniquely determines
\[
  (m_1,\ldots,m_n).
\]
Thus distinct monomials in the $e_i$'s have distinct leading terms.  Consequently, no nontrivial linear combination of monomials in the $e_i$'s can vanish, and there is no nonzero polynomial relation among $e_1,\ldots,e_n$.

\begin{thebibliography}{9}

\bibitem{MS74}
John Milnor and James D. Stasheff,
\emph{Characteristic Classes},
Annals of Mathematics Studies 76,
Princeton University Press, 1974.

\bibitem{Hat02}
Allen Hatcher,
\emph{Algebraic Topology},
Cambridge University Press, 2002.

\end{thebibliography}

\end{document}
